\documentclass[12pt]{article}

\usepackage[a4paper,margin=1in]{geometry}
\usepackage{amsmath, amssymb, booktabs}
\usepackage{setspace}
\usepackage{enumitem}

\setstretch{1.2}

\title{\textbf{Final Exam Review}\\
ECON 204/INTS 211: Contemporary Economic Issues}
\date{}

\begin{document}

\maketitle

%%%%%%%%%%%%%%%%%%%%%%%%%%%%%%%%%%%%%%%%%%%%%%%%%%%%%%
\section*{Part I: Income Distribution and Inequality}

\textbf{Question 1}

Consider the following income distribution (in dollars):

\[
5{,}000,\; 15{,}000,\; 25{,}000,\; 55{,}000,\; 100{,}000
\]

\begin{enumerate}[label=(\alph*)]
\item Compute total income and mean income.
\item Compute the median income.
\item Compute the income share of the top 20\% and bottom 40\%.
\item Suppose the highest income increases to \$140,000. Recompute the mean and comment on inequality.
\item Explain why the mean may be misleading in skewed distributions.
\end{enumerate}

\vspace{1em}

\textbf{Question 2 (Gini Coefficient)}

Using the original income distribution:

\begin{enumerate}[label=(\alph*)]
\item Construct cumulative population shares and cumulative income shares.
\item Using the trapezoidal rule, compute the Gini coefficient:
\[
G = 1 - \sum (Y_i + Y_{i-1})(X_i - X_{i-1})
\]
\item Interpret the magnitude of the Gini coefficient.
\item Explain how a transfer from the richest to the poorest household would affect the Lorenz curve and Gini.
\end{enumerate}

%%%%%%%%%%%%%%%%%%%%%%%%%%%%%%%%%%%%%%%%%%%%%%%%%%%%%%
\section*{Part II: Inflation and Cost of Living}

\textbf{Question 3}

A consumer price index evolves as follows:

\[
CPI_{2024} = 132.5, \quad CPI_{2025} = 140.8
\]

Nominal wages increase from \$52,000 to \$54,600.

\begin{enumerate}[label=(\alph*)]
\item Compute the inflation rate.
\item Compute nominal wage growth.
\item Compute real wage growth.
\item Determine whether purchasing power increased or decreased.
\item Explain how inflation can redistribute income between borrowers and lenders.
\end{enumerate}

\vspace{1em}

\textbf{Question 4 (Household Inflation)}

A household has the following expenditure shares:

\begin{center}
\begin{tabular}{lcc}
\toprule
Category & Weight & Inflation \\
\midrule
Housing & 45\% & 8\% \\
Food & 20\% & 6\% \\
Transport & 15\% & 4\% \\
Other & 20\% & 2\% \\
\bottomrule
\end{tabular}
\end{center}

\begin{enumerate}[label=(\alph*)]
\item Compute the household-specific inflation rate.
\item If official CPI inflation is 3.5\%, explain the discrepancy.
\item Which households are most affected by this inflation pattern?
\end{enumerate}
\newpage
%%%%%%%%%%%%%%%%%%%%%%%%%%%%%%%%%%%%%%%%%%%%%%%%%%%%%%
\section*{Part III: Monetary Policy and Interest Rates}

\textbf{Question 5}

Suppose:

\[
i = 6.5\%, \quad \pi^e = 3.2\%
\]

\begin{enumerate}[label=(\alph*)]
\item Compute the real interest rate.
\item Explain how a rise in the policy rate affects:
\begin{enumerate}
\item consumption
\item investment
\item exchange rates
\end{enumerate}
\item Explain how monetary policy affects inflation expectations.
\end{enumerate}

%%%%%%%%%%%%%%%%%%%%%%%%%%%%%%%%%%%%%%%%%%%%%%%%%%%%%%
\section*{Part IV: Labour Markets and Inequality}

\textbf{Question 6}

Wages for two groups evolve as follows:

\begin{center}
\begin{tabular}{lcc}
\toprule
Group & 2005 & 2025 \\
\midrule
High-skill & \$28 & \$70 \\
Low-skill & \$20 & \$25 \\
\bottomrule
\end{tabular}
\end{center}

\begin{enumerate}[label=(\alph*)]
\item Compute the wage gap in both years.
\item Compute percentage wage growth for each group.
\item Explain how skill-biased technological change explains the results.
\item Discuss one policy that could reduce inequality.
\end{enumerate}
\newpage
%%%%%%%%%%%%%%%%%%%%%%%%%%%%%%%%%%%%%%%%%%%%%%%%%%%%%%
\section*{Part V: Housing Affordability}

\textbf{Question 7}

A household earns \$95,000 annually and considers purchasing a home priced at \$855,000.

\begin{enumerate}[label=(\alph*)]
\item Compute the price-to-income ratio (PTI).
\item Interpret the result.
\item If interest rates increase, explain how affordability changes even if prices remain constant.
\end{enumerate}

\vspace{1em}

\textbf{Question 8 (Mortgage Burden)}

A household takes a mortgage of \$600,000.

\begin{enumerate}[label=(\alph*)]
\item Explain qualitatively how monthly payments change when interest rates rise.
\item Define and explain a \textbf{mortgage renewal shock}.
\item Explain why this is particularly important in Canada.
\end{enumerate}

\vspace{1em}

\textbf{Question 9 (Rental Market)}

Monthly income = \$4,200.

Rent increases from \$1,300 to \$2,050.

\begin{enumerate}[label=(\alph*)]
\item Compute rent-to-income ratios before and after.
\item Explain affordability thresholds.
\item Discuss how rising rents contribute to inequality.
\end{enumerate}
\newpage
%%%%%%%%%%%%%%%%%%%%%%%%%%%%%%%%%%%%%%%%%%%%%%%%%%%%%%
\section*{Part VI: Integrated Policy Analysis}

\textbf{Question 10}

Consider a policy that increases housing supply through zoning reform.

\begin{enumerate}[label=(\alph*)]
\item Explain how this affects housing prices in the short run and long run.
\item Discuss the effect on inequality.
\item Identify one unintended consequence.
\end{enumerate}

\vspace{1em}

\textbf{Question 11 (Integrated Analysis)}

Suppose inflation rises to 6\% and the central bank increases interest rates significantly.

\begin{enumerate}[label=(\alph*)]
\item Explain the impact on:
\begin{enumerate}
\item housing demand
\item mortgage payments
\item rental markets
\end{enumerate}
\item Explain why inflation may fall but housing affordability may worsen.
\item Discuss distributional effects across income groups.
\end{enumerate}

%%%%%%%%%%%%%%%%%%%%%%%%%%%%%%%%%%%%%%%%%%%%%%%%%%%%%%
\section*{End of Review}

\end{document}